\documentclass{situationist-preprint}

\title{The Plurality of Scientific Methods}
\msid{SCI-05-a}
\msfamily{SCI-05}
\msdiscipline{Philosophy of Science}
\mstype{Principal manuscript}
\mskeywords{scientific method; methodological pluralism; Feyerabend; Kuhn; epistemic norms; experiment; historical science; simulation; demarcation}

\begin{document}
\maketitle

\begin{abstract}
Is there a scientific method? The tradition from Bacon to Popper sought a single procedure that every science follows and that explains its success. Feyerabend concluded from the failure of the search that there is no method at all and that, as far as rules go, anything goes. We argue for an intermediate thesis: there is no single scientific method, but there is a small family of methods, controlled experiment, observational inference, modeling and simulation, historical reconstruction and theoretical derivation, united by common norms. We first show, by a dilemma, that no candidate for the single method is both contentful enough to exclude non-science and broad enough to cover the legitimate procedures of the mature sciences. We then identify four norms, answerability to evidence, public checkability, proportioning belief to evidence and the pursuit of severe tests, and argue that every method in the family satisfies them when properly used, that they exclude the procedures of pseudo-science, and that Feyerabend's historical examples violate particular methodological rules but not these norms. The result explains why demarcation does not require a single method, and why scientific education should teach a family of methods rather than one.
\end{abstract}

\section{The problem}

School textbooks describe \enquote{the scientific method} as a sequence of steps: observe, hypothesize, predict, test, conclude. Few working scientists recognize their own practice in the description, and the philosophers who have tried to state the method precisely have disagreed about what it is \citep{Bauer1992}. The question whether there is a scientific method is therefore partly a question about this disagreement: does it reflect the difficulty of stating a single method that exists, or the absence of one?

Three answers dominate. \emph{Methodological monism} holds that there is one method, though its proponents disagree about which: eliminative induction \citep{Bacon1620}, the hypothetico-deductive method \citep{Hempel1966}, falsificationism \citep{Popper1959} or Bayesian conditionalization \citep{HowsonUrbach2006}. \emph{Methodological anarchism} holds that there is no method, that every rule proposed has been violated in the course of some scientific advance, and that the only principle that does not inhibit progress is that anything goes \citep{Feyerabend1975}. \emph{Methodological pluralism} holds that there are many methods, and comes in a strong form, on which the sciences are disunified and share nothing methodologically significant \citep{Dupre1993}, and a weak form, on which the methods share some features that make them scientific \citep{HoyningenHuene2013}.

We defend a version of weak pluralism. There is no single scientific method, but there is a small family of methods united by common norms. The thesis has three parts. The negative part says that no single method is followed by all the sciences. The positive part says that the methods actually used fall into a small number of kinds. The unifying part says that these kinds share a set of norms, which are not themselves a method but constraints on the use of any method, and which distinguish the family from the procedures of non-science.

Section~2 traces the search for the single method, Sections~3 and 4 give the framework and argument, and Sections~5--8 discuss cases, objections, recent literature and limits.

\section{The search for the method}

\paragraph{Bacon and eliminative induction.} \citet{Bacon1620} proposed that natural philosophy proceed by compiling tables of presence, absence and degrees, and by eliminating candidate natures until the form of a phenomenon remained. It failed to describe the sciences of the following century: Newton's gravitation was not arrived at by eliminating natures from tables of heavy things.

\paragraph{Newton and Mill.} Newton's rules of reasoning in Book III of the \emph{Principia} required that no more causes be admitted than are true and sufficient to explain the phenomena, and that propositions inferred by induction from phenomena be held as accurate until other phenomena make them more accurate or liable to exceptions. \citet{Mill1843} codified inductive inference in his methods of agreement, difference, residues and concomitant variation. Mill's methods describe the logic of controlled comparison, and the method of difference remains the logical core of experimental design. But they presuppose that the relevant causal factors have already been identified and can be varied, and they say nothing about how theories of unobservable structure are formed or tested.

\paragraph{Whewell and the hypothetico-deductive method.} \citet{Whewell1840} replied that the sciences proceed by the superinduction of a conception on the facts, a hypothesis that binds them together, and that the mark of a true hypothesis is consilience, its explaining facts of a kind other than those it was introduced to explain. The hypothetico-deductive method, as later formulated by \citet{Hempel1966}, says that a hypothesis is confirmed when observational consequences deduced from it, with auxiliaries, are verified. The method is too permissive: if $H$ entails $E$, so does $H \wedge X$ for any $X$, and the verification of $E$ would confirm the irrelevant conjunction \citep{Glymour1980}. It is also silent on the sciences in which hypotheses are not tested by deducing predictions, such as the historical sciences, whose hypotheses concern unrepeatable past events.

\paragraph{Popper and falsification.} \citet{Popper1959} held that hypotheses are never confirmed, only corroborated by surviving serious attempts at refutation, and that the scientific attitude consists in proposing bold conjectures and trying to falsify them. Falsificationism faced the Duhem problem, since a failed prediction refutes only a conjunction of hypothesis and auxiliaries, and the historical objection that scientists rightly retain theories in the face of anomalies: Newtonian mechanics was retained despite the anomalous perihelion of Mercury for half a century, and rightly so.

\paragraph{Kuhn and Lakatos.} \citet{Kuhn1962} argued that most science is normal science, puzzle-solving within a paradigm that is not itself tested, and that paradigms change by revolutions in which the standards of evaluation themselves shift. \citet{Lakatos1970} tried to reconcile Popper and Kuhn: the unit of appraisal is a research programme, a sequence of theories sharing a hard core, and a programme is progressive if its modifications predict novel facts, degenerating if they merely accommodate. As \citet{Feyerabend1975} pointed out, the methodology gives no rule for when to abandon a degenerating programme, since a programme may recover.

\paragraph{Feyerabend.} \citet{Feyerabend1975} drew the conclusion that no methodological rule is both substantive and universally respected by good science. His central case was Galileo's defence of Copernicanism, which on his reading succeeded through propaganda, ad hoc hypotheses and the use of a telescope whose reliability for celestial observation had not been established, in violation of the methodological rules of his time and of the rules later philosophers proposed. The only rule that survives all such cases, he concluded, is that anything goes.

\paragraph{Laudan.} \citet{Laudan1984} replied that methodological rules are hypothetical imperatives connecting means to cognitive aims, and that they can be evaluated empirically by asking whether following them has in fact promoted those aims. Rules and aims change, but rationally when guided by evidence about which rules work. Laudan's view allows methods to be plural and revisable without being arbitrary.

The pattern is instructive. Each proposed single method captured a real feature of some scientific practice and was defeated by a practice it did not fit. Feyerabend took the pattern to show that there is no method. It shows something weaker: that there is no \emph{single} method, and that the procedures that each candidate captured are among several that science uses.

\section{Framework}

\begin{definition}[Method]\label{d:method}
A \emph{method} is a general procedure that takes evidence of a specified kind and yields conclusions of a specified kind, with conditions under which it is to be applied and standards for its correct application.
\end{definition}

Controlled experiment is a method in this sense: it takes the outcomes of interventions on a system, under conditions in which other causal factors are held fixed or randomized, and yields conclusions about the effects of the intervened variable. Radiometric dating is a method; so is the Monte Carlo simulation of a physical system.

\begin{definition}[Norm]\label{d:norm}
A \emph{norm} is a standard that the use of any method in a domain must satisfy for its conclusions to count as warranted, but that does not itself specify a procedure.
\end{definition}

The distinction between methods and norms is central. A method tells an investigator what to do; a norm tells her what any procedure she uses must achieve. \enquote{Randomize the assignment of treatments} is part of a method; \enquote{do not ignore evidence that bears against your hypothesis} is a norm.

\begin{definition}[Single method thesis]\label{d:smt}
The \emph{single method thesis} holds that there is a method $M$ such that (i) the warranted conclusions of every mature science are reached or justified by $M$, and (ii) $M$ excludes some procedures that are used outside science to reach conclusions about the natural world.
\end{definition}

Clause (ii) gives the thesis content. A method that excluded nothing, such as \enquote{form beliefs about nature somehow}, would trivially satisfy clause (i), and it is not what the defenders of the single method have meant: Bacon, Popper and the textbooks all intend the method to distinguish science from what is not science.

We assume a datum about the mature sciences.

\begin{assumption}[Diversity of procedure]\label{a:div}
The warranted conclusions of the mature sciences are reached by procedures of at least the following kinds: (M1) controlled intervention on a system; (M2) statistical inference from observational data without intervention; (M3) construction and study of models and computer simulations; (M4) reconstruction of past events from their present traces; (M5) mathematical derivation from established theory.
\end{assumption}

Each kind is well established. Experimental physics and much of biology use (M1). Epidemiology, observational astronomy and much of economics use (M2), including the causal inference from correlations and causal graphs that SCI-16 defends. Climate science and much of physics use (M3), whose models are idealized, as SCI-10 argues. Geology, paleontology and evolutionary biology use (M4) \citep{Cleland2002}. Theoretical physics and theoretical population genetics use (M5). Exploratory experimentation, which varies conditions to discover regularities before any theory is available \citep{Steinle1997}, is a variant of (M1).

\section{The argument}

\begin{proposition}\label{p:nosingle}
The single method thesis is false.
\end{proposition}

\begin{proof}
Suppose $M$ is a method satisfying clause (i) of Definition~\ref{d:smt}. By Assumption~\ref{a:div}, $M$ must justify conclusions reached by procedures of all five kinds. These procedures take different kinds of evidence and have different conditions of application: (M1) requires that the investigator can intervene on the system, which is impossible for the past and for most astronomical systems; (M4) requires present traces of past events and draws its strength from the convergence of independent traces, not from intervention; (M3) requires a model whose relation to the target system is assessed by its partial similarity, not by its truth; (M5) uses no new evidence at all. There are then two possibilities.

If $M$ specifies a procedure with conditions of application, as Definition~\ref{d:method} requires, it specifies either a procedure that one of these kinds uses and others do not, in which case it fails clause (i) for the others, or a procedure that all of them share. The only procedures all five share are very general: forming hypotheses, deriving consequences, comparing them with evidence, and revising. A method consisting only of these steps excludes nothing that astrology or numerology does, since the astrologer too forms hypotheses, derives horoscopes and compares them with what happens. Such an $M$ fails clause (ii).

If $M$ does not specify a procedure but only a standard that any procedure must meet, then $M$ is a norm by Definition~\ref{d:norm}, not a method by Definition~\ref{d:method}, and the single method thesis is not satisfied by it. In either case, no method satisfies both clauses.
\end{proof}

The proof uses no premise about the inadequacy of any particular candidate. It shows that the diversity of evidential situations in which the sciences work forces any method that covers them all to be too general to exclude anything, and that what can be general and still exclude is a norm. This is the precise sense in which the search described in Section~2 was bound to fail.

\begin{principle}[Answerability]\label{n:answer}
A warranted conclusion about the natural world must be one that evidence obtained from the world could tell against.
\end{principle}

\begin{principle}[Public checkability]\label{n:public}
The evidence and reasoning for a warranted conclusion must be available to others in a form that allows them to check it, and the conclusion must be exposed to their criticism.
\end{principle}

\begin{principle}[Proportion]\label{n:prop}
Confidence in a conclusion must be proportioned to the evidence, in the sense that it is coherent and responsive to the bearing of the evidence on the conclusion and its rivals.
\end{principle}

\begin{principle}[Severity]\label{n:severe}
The evidence for a conclusion must include the outcome of procedures that would probably have revealed the conclusion to be false if it were false.
\end{principle}

Principle~\ref{n:answer} is Popper's insight, detached from his method. Principle~\ref{n:public} corresponds to the communism and organized skepticism that \citet{Merton1973} identified as norms of science, and to the institutionalized error-correction that SCI-01 takes to distinguish science from non-science; \citet{Longino1990} argues that objectivity in science is secured by this kind of public criticism. Principle~\ref{n:prop} is satisfied by the Bayesian account of confirmation that SCI-12 defends, which states what proportioning confidence to evidence requires. Principle~\ref{n:severe} is the severity requirement of \citet{Mayo1996}, which generalizes the demand for novel prediction.

\begin{proposition}\label{p:satisfy}
Each of the methods (M1)--(M5), properly applied, satisfies Principles~\ref{n:answer}--\ref{n:severe}.
\end{proposition}

\begin{proof}
For (M1), randomized or controlled intervention is a procedure designed to make the observed outcome depend on the effect of the intervened variable, so that the outcome could tell against the hypothesis of an effect (Principle~\ref{n:answer}), and controls and replication are devices for making it probable that a false hypothesis would be exposed (Principle~\ref{n:severe}). For (M2), statistical inference states the probability of the observed data under rival hypotheses and requires that the analysis be specified so that others can reproduce it. For (M3), a simulation is checked against observations of the target system and its robustness to variations in idealizing assumptions is tested, which exposes a model whose success depends on a particular idealization. For (M4), a reconstruction of the past predicts what further traces will be found, and the discovery of independent traces, such as the iridium anomaly and the Chicxulub crater for the impact hypothesis of the end-Cretaceous extinction \citep{Alvarez1980}, is the severe test that \citet{Cleland2002} calls a smoking gun. For (M5), a derivation is checked by others and yields consequences that can be compared with observation. In every case, the method's standards of correct application include publication of evidence and reasoning, so that Principle~\ref{n:public} is satisfied, and its conclusions are held with confidence that varies with the outcome of these checks, so that Principle~\ref{n:prop} is satisfied.
\end{proof}

\begin{proposition}\label{p:exclude}
Procedures that violate one of Principles~\ref{n:answer}--\ref{n:severe} do not yield warranted conclusions, and the characteristic procedures of the standard examples of pseudo-science violate at least one.
\end{proposition}

\begin{proof}
The first part is immediate from Definition~\ref{d:norm}, given that the principles are norms. For the second, astrology, as practised, interprets outcomes so that no outcome tells against its predictions, violating Principle~\ref{n:answer}, and does not subject its predictions to tests that would probably expose them if false, violating Principle~\ref{n:severe}. Creation science holds its central conclusions fixed whatever the evidence, violating Principle~\ref{n:prop}. Practices that rely on private revelation or on the authority of a founder violate Principle~\ref{n:public}.
\end{proof}

\begin{theorem}\label{t:main}
There is no single scientific method, but the methods of the mature sciences form a small family of kinds, (M1)--(M5), united by the norms of Principles~\ref{n:answer}--\ref{n:severe}, which their proper use satisfies and which the procedures of pseudo-science violate.
\end{theorem}

\begin{proof}
The first conjunct is Proposition~\ref{p:nosingle}. The family is given by Assumption~\ref{a:div}; it is small because each kind is individuated by the kind of evidence it takes, and the kinds of evidential situation are few: one can intervene, observe without intervening, model, read traces of the past, or derive. The unity is given by Propositions~\ref{p:satisfy} and \ref{p:exclude}.
\end{proof}

\begin{corollary}\label{c:demarc}
Demarcating science from non-science does not require a single scientific method.
\end{corollary}

\begin{proof}
By Theorem~\ref{t:main}, the norms of Principles~\ref{n:answer}--\ref{n:severe} are shared by the methods of science and violated by the procedures of pseudo-science, and so they draw the distinction that a single method was meant to draw.
\end{proof}

Corollary~\ref{c:demarc} fits the account of demarcation in SCI-01, on which science is distinguished by institutionalized error-correction. Principle~\ref{n:public} is the norm that institutions such as peer review and replication implement, and the other principles describe what the criticism they organize is criticism for.

\section{Cases}

\paragraph{Galileo and the telescope.} Feyerabend's central case is Galileo's use of the telescope to argue for Copernicanism. Feyerabend is right that the reliability of the instrument for celestial objects had not been established by an independent theory of optics, and that Galileo's arguments violated the Aristotelian rule that the senses are reliable only in their natural conditions. But Galileo made his observations public, invited others to repeat them, calibrated the instrument on terrestrial objects whose appearance could be checked, and pointed to the consistency of his observations of Jupiter's satellites over many nights, which no optical illusion was likely to produce. These are applications of Principles~\ref{n:public} and \ref{n:severe}. What Galileo violated was a particular methodological rule, not the norms, and the episode shows that methodological rules may be revised under the guidance of the norms, as \citet{Laudan1984} argued.

\paragraph{The discovery of the Higgs boson.} The discovery combined (M5), the derivation of the boson's production and decay channels from the Standard Model; (M3), the simulation of the expected signal and background in the detectors; and (M1), the controlled collision of protons. The criterion of five standard deviations is a severity requirement, and the use of two independent detectors is an application of Principle~\ref{n:public} within the experiment. No single method describes the episode; the norms describe what made its conclusion warranted.

\paragraph{The impact hypothesis.} The hypothesis that an asteroid impact caused the end-Cretaceous extinction was proposed on the basis of an iridium anomaly at the boundary \citep{Alvarez1980}. It could not be tested by intervention. It was tested by deriving predictions about further traces, shocked quartz, tektites, a crater of the right age and size, and by finding them. The method is (M4), and its success did not depend on resembling the textbook method, which presupposes experiment.

\paragraph{Climate models.} Projections of global temperature under emission scenarios are produced by simulations whose components are derived from physical theory and calibrated against observations \citep{Winsberg2010}. The method is (M3) combined with (M5). Its conclusions are warranted to the extent that the models reproduce observed climate, are robust across independently constructed models and are open to criticism (Principles~\ref{n:severe}, \ref{n:prop} and \ref{n:public}).

\section{Objections}

\begin{objection}[Bayesianism is the method]
The Bayesian account of confirmation applies to every science: in each, credences are updated by conditionalization on evidence. That is a single method, and SCI-12 defends it.
\end{objection}

\begin{reply}
Bayesian conditionalization says how confidence should respond to evidence once hypotheses, evidence and likelihoods are given. It does not say how to generate hypotheses, which evidence to collect, how to design an experiment or a simulation, or how to read a trace of the past, and these are what distinguish the methods of Assumption~\ref{a:div}. In the terms of Definition~\ref{d:norm}, it is a norm, the precise form of Principle~\ref{n:prop}, not a method. SCI-12 is an account of confirmation, and Theorem~\ref{t:main} assigns it exactly that role.
\end{reply}

\begin{objection}[The norms are a method]
\enquote{Formulate testable hypotheses, test them severely, publish the results and proportion belief to the outcome} is a procedure, and it is followed by every science. It is the single method under another name.
\end{objection}

\begin{reply}
If the formula is read as a procedure, it is the general procedure considered in the proof of Proposition~\ref{p:nosingle}, and it specifies no conditions of application or standards of correct application that would allow an investigator to carry it out without already possessing one of the methods (M1)--(M5). To test a hypothesis about the past severely, one must know how to read traces; to test a hypothesis about a drug, one must know how to design a trial. The formula states what these methods must achieve, which is what a norm does.
\end{reply}

\begin{objection}[Anything goes]
Feyerabend showed that every methodological rule has been violated by some scientific advance. The norms are rules too, and they will have been violated.
\end{objection}

\begin{reply}
Feyerabend's examples are violations of methodological rules, such as the rule that new theories must be consistent with well-confirmed old ones, or that instruments must be validated by theory before use \citep{Feyerabend1975}. They are not cases in which a scientific advance was achieved by concealing evidence from criticism, by holding a conclusion fixed whatever the evidence, or by avoiding tests that might refute it. Where such things happened, as in some of Galileo's rhetoric, they are not what made the advance an advance: the Copernican system prevailed because its predictions, and later Kepler's and Newton's, survived tests that the Ptolemaic system failed. Feyerabend's own argument appeals to the norms when it claims that Galileo's theory was in the end better supported.
\end{reply}

\begin{objection}[Paradigm relativity]
\citet{Kuhn1962} argued that the standards of evaluation change with paradigms. The norms are then paradigm-relative, and the unity they provide is illusory.
\end{objection}

\begin{reply}
What changes across revolutions, on Kuhn's own later account, is the weighting and interpretation of shared values, such as accuracy, consistency, scope, simplicity and fruitfulness, not the values themselves \citep{Kuhn1977}. Principles~\ref{n:answer}--\ref{n:severe} are at the level of these shared values, and they are respected on both sides of the revolutions Kuhn describes. SCI-06 argues that successive paradigms are partially translatable and that theory choice across revolutions remains rational; the shared norms are part of what makes it so.
\end{reply}

\begin{objection}[Disunity]
\citet{Dupre1993} argues that the sciences are disunified in their ontology and their methods, and that any norms they share are too thin to constitute a unity. A family united only by thin norms is no family.
\end{objection}

\begin{reply}
The norms are thin in that they specify no procedure, but they are not empty: Proposition~\ref{p:exclude} shows that they exclude the procedures of pseudo-science, and Proposition~\ref{p:satisfy} shows that they are what the specific standards of each method serve. Unity of this kind is what is needed to explain why the methods are all scientific and why their conclusions can be combined, as in the Higgs case, where three methods were used together and their results were assessed by common standards. Dupré's disunity of ontology is compatible with this unity of norms.
\end{reply}

\section{Recent literature and consequences}

\citet{HoyningenHuene2013} argues that what distinguishes scientific knowledge is not a method but its greater systematicity, along dimensions such as description, explanation, prediction, defence of knowledge claims and critical discourse. His view and ours agree that there is no single method and that science is unified by something more abstract. They differ in that systematicity is a matter of degree along many dimensions, which makes it hard to say why astrology, which is highly systematic in its own way, is not science. Our norms exclude astrology by Proposition~\ref{p:exclude}. \citet{Crombie1994} distinguished six styles of scientific thinking in the European tradition, including postulation in mathematics, experimental exploration, hypothetical modeling, taxonomy, statistical analysis and historical derivation, and \citet{Hacking1992} developed the idea into an account of styles of reasoning that bring their own objects and criteria of truth into being. Crombie's styles correspond closely to the methods of Assumption~\ref{a:div}, with taxonomy as a further candidate, and support the thesis that the family is small. We do not follow Hacking in holding that each style constitutes its own criteria of truth; the norms are common to the styles.

Two consequences follow for practice. First, science education that presents a single method misdescribes science and leaves students unable to recognize the historical and modeling sciences as sciences \citep{Bauer1992}. It should teach a family of methods and the norms that unite them. Second, public assessment of scientific claims should ask whether the norms were satisfied, not whether the claim was reached by the textbook method. Claims about climate, reached by modeling, and about evolutionary history, reached by reconstruction, are sometimes dismissed as unscientific because they were not established by experiment; the dismissal rests on the single method thesis, which is false.

\section{What the argument does not establish}

\begin{remark}\label{r:limits}
The argument does not establish that the list (M1)--(M5) is complete or that its boundaries are sharp: taxonomy and exploratory experimentation may deserve separate places, and many studies combine methods. It does not establish that the four norms are the only common norms; consistency and the avoidance of fraud are plausible further candidates, though the latter is arguably implied by Principle~\ref{n:public}. And it does not establish how the norms are to be weighed when they conflict, for instance when severe testing is impossible and a conclusion must be drawn from a single unrepeatable observation. These are matters for the specific methods and for the judgement of investigators, which the norms constrain but do not replace.
\end{remark}

\section{Conclusion}

There is no single scientific method. Any procedure general enough to cover experiment, observational inference, modeling, historical reconstruction and theoretical derivation is too general to exclude the procedures of pseudo-science, and what is general enough to cover them and still excludes is a norm, not a method. The methods of science form a small family, individuated by the evidential situations in which science works, and united by four norms: answerability to evidence, public checkability, proportioning confidence to evidence and severe testing. Feyerabend was right that every methodological rule has been broken in good science and wrong that anything goes; Kuhn was right that standards are interpreted differently across paradigms and wrong, on his own later account, that they are wholly replaced.

\begin{thebibliography}{99}\small
\bibitem[Alvarez et~al.(1980)]{Alvarez1980} Alvarez, L.~W., Alvarez, W., Asaro, F. and Michel, H.~V. (1980). Extraterrestrial cause for the Cretaceous--Tertiary extinction. \emph{Science}, 208(4448), 1095--1108.
\bibitem[Bacon(1620)]{Bacon1620} Bacon, F. (1620). \emph{Novum Organum}. Book II.
\bibitem[Bauer(1992)]{Bauer1992} Bauer, H.~H. (1992). \emph{Scientific Literacy and the Myth of the Scientific Method}. University of Illinois Press.
\bibitem[Cleland(2002)]{Cleland2002} Cleland, C.~E. (2002). Methodological and epistemic differences between historical science and experimental science. \emph{Philosophy of Science}, 69(3), 474--496.
\bibitem[Crombie(1994)]{Crombie1994} Crombie, A.~C. (1994). \emph{Styles of Scientific Thinking in the European Tradition}. Duckworth.
\bibitem[Dupr\'e(1993)]{Dupre1993} Dupr\'e, J. (1993). \emph{The Disorder of Things: Metaphysical Foundations of the Disunity of Science}. Harvard University Press.
\bibitem[Feyerabend(1975)]{Feyerabend1975} Feyerabend, P. (1975). \emph{Against Method}. New Left Books.
\bibitem[Glymour(1980)]{Glymour1980} Glymour, C. (1980). \emph{Theory and Evidence}. Princeton University Press.
\bibitem[Hacking(1992)]{Hacking1992} Hacking, I. (1992). \enquote{Style} for historians and philosophers. \emph{Studies in History and Philosophy of Science}, 23(1), 1--20.
\bibitem[Hempel(1966)]{Hempel1966} Hempel, C.~G. (1966). \emph{Philosophy of Natural Science}. Prentice-Hall.
\bibitem[Howson and Urbach(2006)]{HowsonUrbach2006} Howson, C. and Urbach, P. (2006). \emph{Scientific Reasoning: The Bayesian Approach}. 3rd edition. Open Court.
\bibitem[Hoyningen-Huene(2013)]{HoyningenHuene2013} Hoyningen-Huene, P. (2013). \emph{Systematicity: The Nature of Science}. Oxford University Press.
\bibitem[Kuhn(1962)]{Kuhn1962} Kuhn, T.~S. (1962). \emph{The Structure of Scientific Revolutions}. University of Chicago Press.
\bibitem[Kuhn(1977)]{Kuhn1977} Kuhn, T.~S. (1977). Objectivity, value judgment, and theory choice. In \emph{The Essential Tension}. University of Chicago Press.
\bibitem[Lakatos(1970)]{Lakatos1970} Lakatos, I. (1970). Falsification and the methodology of scientific research programmes. In I.~Lakatos and A.~Musgrave (eds.), \emph{Criticism and the Growth of Knowledge}. Cambridge University Press.
\bibitem[Laudan(1984)]{Laudan1984} Laudan, L. (1984). \emph{Science and Values}. University of California Press.
\bibitem[Longino(1990)]{Longino1990} Longino, H.~E. (1990). \emph{Science as Social Knowledge}. Princeton University Press.
\bibitem[Mayo(1996)]{Mayo1996} Mayo, D.~G. (1996). \emph{Error and the Growth of Experimental Knowledge}. University of Chicago Press.
\bibitem[Merton(1973)]{Merton1973} Merton, R.~K. (1973). \emph{The Sociology of Science: Theoretical and Empirical Investigations}. University of Chicago Press.
\bibitem[Mill(1843)]{Mill1843} Mill, J.~S. (1843). \emph{A System of Logic}. John W. Parker.
\bibitem[Popper(1959)]{Popper1959} Popper, K.~R. (1959). \emph{The Logic of Scientific Discovery}. Hutchinson.
\bibitem[Steinle(1997)]{Steinle1997} Steinle, F. (1997). Entering new fields: Exploratory uses of experimentation. \emph{Philosophy of Science}, 64 (Proceedings), S65--S74.
\bibitem[Whewell(1840)]{Whewell1840} Whewell, W. (1840). \emph{The Philosophy of the Inductive Sciences}. John W. Parker.
\bibitem[Winsberg(2010)]{Winsberg2010} Winsberg, E. (2010). \emph{Science in the Age of Computer Simulation}. University of Chicago Press.
\end{thebibliography}

\end{document}
